\documentclass[aspectratio=169]{beamer}

\usetheme{Madrid}
\useoutertheme[subsection=false]{miniframes}

\usepackage{amsmath}
\usepackage{amssymb}
\usepackage{mathtools}
\usepackage{graphicx}
\usepackage{braket}
\usepackage{caption}
\usepackage{subcaption}
\usepackage[backend=biber,style=authoryear]{biblatex}

\addbibresource{references.bib}

\setbeamertemplate{navigation symbols}{}

\title[PXP Polynomial Filtering]
{Targeting Scar States Using Polynomial Filtering in the PXP Hamiltonian}

\author{Ahmet Berk Kavruk}
\date{July 3, 2026}

\begin{document}

\begin{frame}[plain,noframenumbering]
    \titlepage
\end{frame}

\begin{frame}{Eigenstate Thermalization Hypothesis (ETH)}
    \begin{itemize}
        \item Explains thermalization in isolated quantum many-body systems.
        \item Local observables relax to match expected values from statistical mechanics.
        \item Individual generic eigenstates mimic thermal ensembles.
    \end{itemize}
    \vfill
    \begin{block}{Important: ETH Violations}
        Some systems have rare, unusual eigenstates that strongly violate ETH and fail to thermalize.
    \end{block}
\end{frame}

\begin{frame}{The PXP Model \& Quantum Many-Body Scars}
    \begin{itemize}
        \item Models projected spin-flip dynamics on a 1D Rydberg atom chain:
        \begin{equation}
            H = \sum_{i=0}^{L-1} P_{i} X_{(i+1)\%L} P_{(i+2)\%L}
        \end{equation}
        where $P_i = \frac{1}{2}(\mathbb{1} - Z_i)$ projects site $i$ onto its ground state $\ket{0}$.
        \item ETH-violating eigenstates form \textbf{quantum many-body scar states}.
        \item Scars are rare, evenly spaced, and possess unusually low entanglement.
        \item But there is the usual problem, the Hilbert space grows exponentially.

        \begin{block}{Problem of the Project}
            Can we use Chebyshev polynomial filtering to target the scar states without doing full exact diagonalization?
        \end{block}
    \end{itemize}
\end{frame}

\begin{frame}{Understanding the PXP Hamiltonian}
    Adjacency graph etc.
\end{frame}


\end{document}